\section{Prompting en tiempo de inferencia y presupuesto de tokens}

\subsection{Ampliar el presupuesto de tokens para exponer la cadena de razonamiento}
Chen usa el rompecabezas de Arc AGI como ejemplo para mostrar la curva del presupuesto de tokens en tiempo de inferencia: los modelos tradicionales se detienen por debajo del 25\% con un budget de aproximadamente 20 dólares, mientras que OpenAI O3, con un presupuesto de entre 20 y 1,000 dólares y acompañado de prompts de ejemplo y un trace de ejecución, puede llevar la exactitud hasta el 87.5\%. Esta diapositiva (figura \ref{fig:token-budget}) muestra con claridad la tendencia no lineal en la parte media de la curva de precisión de razonamiento a medida que se alarga el presupuesto de tokens.

\begin{figure}[H]
\centering
\includegraphics[width=0.92\textwidth]{frame-token-budget.jpg}
\caption{La curva en tiempo de inferencia de Arc AGI y la disposición de costo-token de OpenAI O3.\protect\footnotemark}
\label{fig:token-budget}
\end{figure}
\footnotetext{Intervalo de tiempo en el video: 00:08:56--00:10:20, usa la curva de O3 para mostrar que la expansión de tokens puede alcanzar una exactitud a nivel humano.}

\begin{importantbox}{Tres lecciones de la curva en tiempo de inferencia}
1) Cuando la complejidad de la tarea aumenta, el presupuesto de tokens debe pasar de decenas a miles para ver una mejora significativa; 2) aumentar el presupuesto debe acompañarse de prompts de ejemplo y de una puntuación por pasos (stepwise scoring), de lo contrario un prompt más largo por sí solo no eleva la exactitud; 3) el modo de presupuesto alto es adecuado para tareas de nivel de competencia o de verificación, no para conversaciones nativas de baja latencia.
\end{importantbox}

\subsection{Restricciones del diseño de prompts}
El prompt necesita dos capacidades: primero, alentar al modelo a generar una Chain-of-Thought larga, para que surja de manera natural una secuencia de razonamiento coherente a lo largo de varias rondas; segundo, incorporar en el prompt una señal de dificultad de la tarea, de modo que el modelo ajuste automáticamente la profundidad de generación. En el minuto 44:31, Chen señala explícitamente: «el prompt solo puede producir un razonamiento más largo si el modelo al que se le da el prompt sabe que enfrenta un problema difícil».

\begin{knowledgebox}{Dos criterios del diseño de prompts}
1. CoT largo: usar plantillas por pasos, pares de pregunta y respuesta o demostraciones con varios ejemplos, para que el modelo dé seguimiento continuo a su estado actual; 2. marca de dificultad: mostrar en el prompt la dificultad del problema (problem difficulty), la longitud de contexto (context length) y la profundidad esperada (expected depth), para que el modelo ajuste de manera adaptativa la profundidad en tiempo de inferencia.
\end{knowledgebox}

\subsection{Observabilidad de costos y control de presupuesto}
Ampliar el presupuesto de tokens debe tomar como referencia el costo real: la diferencia entre 1 dólar, 20 dólares y 1,000 dólares por pregunta no es solo cuestión de dinero, también implica latencia y consumo de cómputo. El riesgo real expuesto es que «un presupuesto alto vuelve la tarea inexplicable y también puede consumir los recursos de los microservicios».

\begin{warningbox}{Riesgo de que el presupuesto se descontrole}
En cuanto el presupuesto en tiempo de inferencia se eleva a miles de dólares, es necesario limitarlo a subtareas de alto valor; de lo contrario, la frecuencia de llamadas queda fuera de sincronía, el rendimiento del sistema cae e incluso puede provocar una pérdida de equilibrio. Se recomienda definir de antemano una estrategia de caché de muestreo (sampling) y reservar la búsqueda (search) más costosa para la ruta crítica verificada (verified-critical path).
\end{warningbox}

\subsection{Resumen de la sección}
El prompting en tiempo de inferencia necesita encontrar un equilibrio entre «alentar un CoT largo» y «percibir la dificultad de la tarea», al mismo tiempo que se incorporan a la supervisión el presupuesto de tokens, el costo y las métricas de latencia. La gráfica de la curva de Arc AGI nos recuerda que ampliar la cantidad de tokens es eficaz, pero el costo se satura de manera gradual, por lo que el prompt, la verificación y el presupuesto deben operar de forma coordinada.

